\chapter{从模形式到自守表示 II}

% ===== PDF p.024 =====
\textbf{今天：}
\begin{itemize}
\item 继续构造与尖点形式相伴的自守表示。
\end{itemize}

\textbf{上次：}
\begin{itemize}
\item 我们构造了 $\GL_2(\Af)$-等变嵌入
\[
\Phi:M_k\to C^\infty(\GL_2(\Q)\backslash\GL_2(\A)),
\]
它诱导嵌入
\[
\widetilde\Phi:S_k\to L^2([\GL_2]),
\]
后者在相差因子 $|\det|^{k/2}_{\mathrm{ad\acute{e}lique}}$ 的意义下是 $\GL_2(\Af)$-等变的；也就是说，态射
\[
\widetilde\Phi:S_k\otimes|\det|^{k/2}_{\mathrm{ad\acute{e}lique}}\to L^2([\GL_2])
\]
是 $\GL_2(\Af)$-等变的。
\item 这里
\[
M_k=\varinjlim_{K\subseteq\GL_2(\Af)}M_k(K),
\]
而
\[
M_k(K)=\bigoplus_{i=1}^mM_k(\Gamma_i)
\]
是同余子群 $\Gamma_i\subseteq\GL_2(\Z)$ 的模形式空间之和。
\item 类似地：
\[
S_k=\varinjlim_KS_k(K),
\]
其中 $S_k(K)$ 是同余子群的尖点形式空间之和。
\item 今天：如何刻画像
\[
\Phi(M_k)\subseteq C^\infty(\GL_2(\Q)\backslash\GL_2(\A))
\]
以及
\[
\widetilde\Phi(S_k)\subseteq L^2([\GL_2])？
\]
\item 下次：$S_k$ 如何作为 $\GL_2(\Af)$-表示分解？
\end{itemize}

上一讲中 $\Phi$ 的构造：
\begin{itemize}
\item 我们考虑了 $\GL_2(\R)$-等变线丛
\[
L_k\to\HH^\pm\cong\GL_2(\R)/\C^\times,
\]
它们相伴于 $\C^\times=A_{\GL_2}\SO_2(\R)$ 的表示 $z\mapsto z^k$。
\item 通过拉回，得到 $\GL_2(\Q)\backslash(\GL_2(\Af)\times\HH^\pm)$ 上的线丛 $\omega^{\otimes k}$。
\item $\omega^{\otimes k}$ 的光滑、全纯……截面等同于满足
\[
f(\gamma g,\gamma z)=j(\gamma,z)^kf(g,z)
\]
的光滑、全纯……函数
\[
f:\GL_2(\Af)\times\HH^\pm\to\C,\qquad(g,z)\mapsto f(g,z).
\]
% ===== PDF p.025 =====
\item 这里
\[
\gamma=\begin{pmatrix}a&b\\c&d\end{pmatrix}\in\GL_2(\Q),\qquad j(\gamma,z):=cz+d.
\]
\item 函数
\[
j:\GL_2(\R)\times\HH^\pm\to\C^\times,\qquad
\Bigl(\begin{pmatrix}a&b\\c&d\end{pmatrix},z\Bigr)\mapsto cz+d
\]
满足「余圈方程」
\[
j(gh,z)=j(g,hz)j(h,z)
\]
对一切 $g,h\in\GL_2(\R)$、$z\in\HH^\pm$ 成立。
\item 对如上 $f:\GL_2(\Af)\times\HH^\pm\to\C$，函数
\[
\phi_f:\GL_2(\Af)\times\GL_2(\R)\to\C,\qquad
(g,g_\infty)\mapsto j(g_\infty,i)^{-k}f(g,g_\infty i)
\]
是 $\GL_2(\Q)$-等变的，即属于 $C^\infty(\GL_2(\Q)\backslash\GL_2(\A))$。
\item 由引理 3.2 得到：映射 $f\mapsto\phi_f$ 把光滑截面
\[
\GL_2(\Q)\backslash(\GL_2(\Af)\times\HH^\pm)\to\omega^{\otimes k}
\]
与满足
\[
\phi(g,g_\infty z)=z^{-k}\phi(g,g_\infty)
\]
（对一切 $(g,g_\infty)\in\GL_2(\Af)\times\GL_2(\R)$、$z\in\C^\times$）的光滑函数
\[
\phi:\GL_2(\Q)\backslash\GL_2(\A)\to\C
\]
等同起来。
\end{itemize}

\textbf{要点与下一步目标：}
\begin{itemize}
\item 用 $\phi\in C^\infty(\GL_2(\Q)\backslash\GL_2(\A))$ 上的条件，描述 $M_k$ 在 $f\mapsto\phi_f$ 下的像。
\item 再加上尖点性条件，就决定了 $S_k$ 在
\[
\widetilde\Phi:S_k\to L^2([\GL_2])
\]
下的像：它是那些光滑尖点函数 $\phi$，使得 $|\det|^{-k/2}_{\mathrm{ad\acute{e}lique}}\phi$ 落在 $\Phi$ 的像中。
\item 我们已经知道，$\Phi$ 的像中的元素 $\phi$ 必须满足
\[
\phi(g,g_\infty z)=z^{-k}\phi(g,g_\infty)
\]
（$z\in\C^\times$，$(g,g_\infty)\in\GL_2(\A)$），并且这样的 $\phi$ 下降为满足 $\GL_2(\Q)$-模性的光滑函数
\[
f_\phi:\GL_2(\Af)\times\HH^\pm\to\C,\qquad(g,z)\mapsto f(g,z).
\]
\item 需要检验：
\begin{itemize}
\item $f_\phi$ 何时全纯？
\item $f_\phi$ 何时在尖点处全纯？
\end{itemize}
\end{itemize}

% ===== PDF p.026 =====
\textbf{$f_\phi$ 何时全纯？}这归约到如下问题：给定光滑函数
\[
f:\HH^\pm\to\C,
\]
$\phi_f:\GL_2(\R)\to\C$，$g\mapsto j(g,i)^{-k}f(gi)$ 上的什么条件保证 $f$ 全纯？我们概述将遵循的策略：
\begin{itemize}
\item Lie 代数
\[
\g:=\mathfrak{gl}_2(\R)\cong\mathrm{Mat}_{2,2}(\R)
\]
经右平移对 $\GL_2(\R)$ 的作用求导，而作用在 $C^\infty(\GL_2(\R))$ 上\znote{直观上，$f$ 全纯等价于 $\partial f/\partial\bar z=0$；而 $\partial/\partial\bar z$ 是 $\GL_2(\R)$ 的 Lie 代数作用的一个组合（见下文对 $Y$ 的构造），于是全纯性可以写成 Lie 代数层面的一个方程。这正是把「模形式的全纯性」翻译为「自守表示的无穷小条件」的出发点。[来源：原文第 26--29 页]}。
\item 把这个作用在 $\phi_f$ 上显式化。
\item 然后构造元素 $Y\in\g_\C$，使得
\[
Y*\phi_f=0
\]
当且仅当 $f$ 全纯。这里 $\g_\C:=\g\otimes_\R\C$ 是 $\g$ 的复化 Lie 代数。
\end{itemize}

\textbf{无穷小作用}（参见 \texttt{[GH19, Section 4.2]}）：我们开始解释「对 $\GL_2(\R)$ 的作用求导」是什么意思。
\begin{itemize}
\item 设 $H$ 为任意 Lie 群，Lie 代数 $\mathfrak{h}:=T_1H$。这里 $T_1H$ 表示实流形 $H$ 在单位元处的切空间，即 $\mathfrak{h}=T_1H$ 的元素是道路
\[
\gamma:(-\epsilon,\epsilon)\to H
\]
（$\epsilon>0$，$\gamma(0)=1$）的等价类，两条道路等价若它们在 $1$ 处（在某坐标卡中）有相同的导数。
\item 设 $V$ 是 $H$ 的任意表示，其中 $V$ 是 Hausdorff 拓扑 $\C$-向量空间。
\item 对 $X\in\mathfrak{h}$（由某条道路 $\gamma:(-\epsilon,\epsilon)\to H$，$\gamma(0)=1$ 定义）与 $v\in V$，考虑函数
\[
(-\epsilon,\epsilon)\setminus\{0\}\to V,\qquad t\mapsto\frac{\gamma(t)v-v}{t}.
\]
若该函数在 $t\to0$ 时的极限存在，我们记
\[
X*v:=\lim_{t\to0}\frac{\gamma(t)v-v}{t}.
\]
\item 元素 $v\in V$ 称为\textbf{光滑的}，若
\[
X_1*(\cdots*(X_n*v)\cdots)
\]
对一切 $X_1,\dots,X_n\in\mathfrak{h}$ 都存在。
\item 记 $V^{\mathrm{sm}}\subseteq V$ 为 $V$ 中光滑向量的子空间。
\item 则 $V^{\mathrm{sm}}$ 在 $H$ 下稳定（但可能与 $V$ 不同）。
\item 此外，$V^{\mathrm{sm}}$ 是 $\mathfrak{h}$ 的表示，或等价地是 $\mathfrak{h}$ 的包络代数
\[
U(\mathfrak{h}):=\bigoplus_{n\ge0}\mathfrak{h}^{\otimes n}/\langle X\otimes Y-Y\otimes X-[X,Y]\mid X,Y\in\mathfrak{h}\rangle
\]
上的模。
\item 这里 $[-,-]$ 是 $\mathfrak{h}$ 的 Lie 括号。
\end{itemize}

% ===== PDF p.027 =====
\textbf{$\g=\mathfrak{gl}_2(\R)$ 在 $\phi_f$ 上的显式作用：}现在我们把 $\mathfrak{gl}_2(\R)$ 在函数
\[
\phi_f(g):=j(g,i)^{-k}f(gi)
\]
上的作用写得更显式，其中 $\phi_f$ 相伴于某个截面
\[
f:\HH^\pm\to\OO_{\HH^\pm}(k).
\]
为此回忆等式
\[
j(gh,z)=j(g,hz)j(h,z)
\]
（$g,h\in\GL_2(\R)$，$z\in\HH^\pm$）。
\begin{itemize}
\item 我们已经知道
\[
\phi_f(gz)=z^{-k}\phi_f(g)
\]
（$z\in\C^\times$，$g\in\GL_2(\R)$，由引理 3.2）。注意
\[
\C\cong\Lie(\C^\times)=\Bigl\langle\begin{pmatrix}1&0\\0&1\end{pmatrix},\begin{pmatrix}0&-1\\1&0\end{pmatrix}\Bigr\rangle_\R\subseteq\g.
\]
求导可推出 $\C$ 通过线性形式
\[
\C\to\C,\qquad z\mapsto-kz
\]
作用在 $\phi_f$ 上。因此更精确地，
\[
\begin{pmatrix}1&0\\0&1\end{pmatrix}*\phi_f=-k\phi_f,
\qquad
\begin{pmatrix}0&-1\\1&0\end{pmatrix}*\phi_f=-ki\phi_f,
\]
因为 $\begin{pmatrix}0&-1\\1&0\end{pmatrix}$ 对应 $i\in\C$。
\item 现在考虑子群
\[
U:=\Bigl\{\begin{pmatrix}*&*\\0&1\end{pmatrix}\Bigr\}\subseteq\GL_2(\R)
\]
及其 Lie 代数
\[
\Lie(U)=\Bigl\langle\begin{pmatrix}1&0\\0&0\end{pmatrix},\begin{pmatrix}0&1\\0&0\end{pmatrix}\Bigr\rangle_\R\subseteq\g.
\]
\item 定义辅助函数
\[
\phi_f^\sharp:\GL_2(\R)\times\HH^\pm\to\C,\qquad(g,z)\mapsto j(g,z)^{-k}f(gz)
\]
（不要把它与上一讲引入的 $\widetilde\phi_f$ 混淆）\znote{原文把这里的辅助函数 $\phi_f^\sharp(g,z):=j(g,z)^{-k}f(gz)$ 与上一讲的 $\widetilde\phi_f=|\det|^{k/2}_{\mathrm{ad\acute{e}lique}}\phi_f(x)$ 用了同一个记号 $\widetilde\phi_f$（原文自己也提醒「不要混淆」）。为免混淆，译文把本讲的这个辅助函数记作 $\phi_f^\sharp$，而上一讲的仍记作 $\widetilde\phi_f$。[来源：原文第 27 页；另见第 3 讲第 21 页]}。
\item 对固定的 $g\in\GL_2(\R)$ 得到
\[
\frac{\partial\phi_f^\sharp}{\partial\bar z}(g,z)=\det(g)j(g,z)^{-k-2}\frac{\partial f}{\partial\bar z}(gz)
\]
（用到了函数 $z\mapsto gz$ 的 $\partial/\partial z$-导数为 $\det(g)/j(g,z)^2$）。
\item 从而 $f$ 全纯当且仅当
\[
\frac{\partial\phi_f^\sharp}{\partial\bar z}(g,i)=0
\]
对一切 $g\in\GL_2(\R)$ 成立。
\end{itemize}
% ===== PDF p.028 =====
\begin{itemize}
\item 对任意道路 $\gamma:(-\epsilon,\epsilon)\to U$，$\gamma(0)=1$，可以算出
\begin{align*}
\phi_f(g\gamma(t))
&=\phi_f^\sharp(g\gamma(t),i)\\
&=j(g\gamma(t),i)^{-k}f(g\gamma(t)i)\\
&=j(g,\gamma(t)i)^{-k}j(\gamma(t),i)^{-k}f(g\gamma(t)i)\\
&=\phi_f^\sharp(g,\gamma(t)i),
\end{align*}
因为 $j(\gamma(t),i)=1$。
\item 令
\[
X_\gamma:=\frac{\partial\gamma}{\partial t}(0)\in\Lie(U).
\]
\item 则
\begin{align*}
X_\gamma*\phi_f
&=\frac{\partial}{\partial t}\Bigl(\phi_f(g\gamma(t))\Bigr)\Big|_{t=0}\\
&=\frac{\partial\phi_f^\sharp}{\partial z}(g,i)\cdot\frac{\partial}{\partial t}\bigl(\gamma(t)\cdot i\bigr)\Big|_{t=0}
+\frac{\partial\phi_f^\sharp}{\partial\bar z}(g,i)\cdot\frac{\partial}{\partial t}\bigl(\overline{\gamma(t)\cdot i}\bigr)\Big|_{t=0},
\end{align*}
这里用了 Wirtinger 导数的链式法则\znote{原文此处（以及第 27、29 页）印刷为 $\partial/\partial z$，但按链式法则第二项应为 $\partial/\partial\bar z$，且后文「$Y*\phi_f=2i\partial\phi_f^\sharp/\partial\bar z$」与「$f$ 全纯 $\iff$ $Y*\phi_f=0$」只有在 $\partial/\partial\bar z$ 时才自洽（提取文本系统性地丢了 $\bar z$ 上的横线，参见第 13 页 $dx\wedge dy=\frac i2\,dz\wedge d\bar z$）。译文按 $\partial/\partial\bar z$ 处理。[来源：原文第 27--29 页]}。
\item 更具体地，先取第一条道路
\[
\gamma_1:(-\epsilon,\epsilon)\to U,\qquad t\mapsto\begin{pmatrix}1+t&0\\0&1\end{pmatrix}.
\]
则 $\gamma_1(t)\cdot i=(1+t)i$，于是
\[
\begin{pmatrix}1&0\\0&0\end{pmatrix}*\phi_f(g)
=i\frac{\partial\phi_f^\sharp}{\partial z}(g,i)-i\frac{\partial\phi_f^\sharp}{\partial\bar z}(g,i).
\]
\item 再取第二条道路
\[
\gamma_2:(-\epsilon,\epsilon)\to U,\qquad t\mapsto\begin{pmatrix}1&t\\0&1\end{pmatrix},
\]
则 $\gamma_2(t)\cdot i=i+t$\znote{原文此处误作 $\gamma_1(t)\cdot i=i+t$；由定义 $\gamma_1(t)\cdot i=(1+t)i$，而 $\gamma_2(t)\cdot i=i+t$，故应为 $\gamma_2$。[来源：原文第 28 页]}，于是
\[
\begin{pmatrix}0&1\\0&0\end{pmatrix}*\phi_f(g)
=\frac{\partial\phi_f^\sharp}{\partial z}(g,i)+\frac{\partial\phi_f^\sharp}{\partial\bar z}(g,i).
\]
\item 特别地，
\[
i\begin{pmatrix}1&0\\0&0\end{pmatrix}*\phi_f(g)+\begin{pmatrix}0&1\\0&0\end{pmatrix}*\phi_f(g)
=2\frac{\partial\phi_f^\sharp}{\partial\bar z}(g,i).
\]
\item 因此
\[
\begin{pmatrix}i&1\\0&0\end{pmatrix}\in\g_\C
\]
可以作为 $Y$ 的候选，即 $f$ 全纯当且仅当
\[
\begin{pmatrix}i&1\\0&0\end{pmatrix}*\phi_f=0.
\]
\item 但我们还可以做得更好。也就是说，定义
\[
Y:=\frac12\begin{pmatrix}-1&i\\i&1\end{pmatrix}
=\frac12\begin{pmatrix}1&0\\0&1\end{pmatrix}+\frac{1}{2i}\begin{pmatrix}0&-1\\1&0\end{pmatrix}
+i\begin{pmatrix}i&1\\0&0\end{pmatrix}.
\]
\znote{原文这个分解式与左边的矩阵不符：按右边逐项相加得到的是 $\frac12\begin{pmatrix}-1&i\\-i&1\end{pmatrix}$，即 $(2,1)$-项的符号相反（把 $+\frac{1}{2i}$ 改为 $-\frac{1}{2i}$ 才与左边的矩阵一致；下页对 $Y*\phi_f$ 的计算实际用的也是与左边矩阵一致的那个元素）。译文照左边的矩阵与原文的分解式直译，仅在此说明该分解式有误，不另作改动。[来源：原文第 28--29 页；参见 work/疑点清单.tsv 第 28 页条目]}
\end{itemize}

% ===== PDF p.029 =====
于是
\[
Y*\phi_f(g)
=\underbrace{-\frac k2\phi_f(g)+i\Bigl(-\frac{ki}{2}\Bigr)\phi_f(g)}_{=0}
+2i\frac{\partial\phi_f^\sharp}{\partial\bar z}(g,i)
=2i\frac{\partial\phi_f^\sharp}{\partial\bar z}(g,i),
\]
从而 $f$ 全纯当且仅当 $Y*\phi_f=0$。
\begin{itemize}
\item 为什么这个选择更好？
\begin{itemize}
\item 令
\[
H:=\begin{pmatrix}0&-i\\i&0\end{pmatrix},\qquad
X:=\frac12\begin{pmatrix}-1&-i\\-i&1\end{pmatrix}.
\]
\item 则 $H,X,Y$ 是 $\g_\C$ 中的一个 $\mathfrak{sl}_2$-三元组。
\item 确切地说，在基
\[
\begin{pmatrix}-i\\1\end{pmatrix},\qquad\begin{pmatrix}i\\1\end{pmatrix}
\]
下，它由
\[
\begin{pmatrix}1&0\\0&-1\end{pmatrix},\qquad
\begin{pmatrix}0&1\\0&0\end{pmatrix},\qquad
\begin{pmatrix}0&0\\1&0\end{pmatrix}
\]
给出。
\item 注意 $H$ 张成 Lie 代数 $i\Lie(\SO_2(\R))\subseteq\g_\C$。
\end{itemize}
\end{itemize}

我们可以转向下一个问题了。

\textbf{$f_\phi$ 何时在尖点处全纯？}
\begin{itemize}
\item 考虑具有 Fourier 展开的函数
\[
f(z)=\sum_{n\in\Z}a_ne^{2\pi izn},
\]
它定义在
\[
\{z\in\C\mid|\Re(z)|<1,\ \Im(z)>1\}
\]
上。
\item 回忆
\[
|e^{2\pi iz}|=e^{-2\pi y},
\]
其中 $y=\Im(z)$。
\item 则 $a_n=0$ 对一切 $n<0$ 成立，当且仅当
\[
|f(x+iy)|\le Cy^N
\]
对某 $C,N\in\R_{>0}$ 与 $y\to\infty$ 成立，因为当 $n<0$ 时指数函数 $e^{-2\pi ny}$ 的增长快于任何多项式函数。
\item 按定义，这就是说 $f$ 具有\textbf{适度增长}（或缓慢增长）；这一概念可以推广到 $\Spec(\Q)$ 上的一切约化群。
\item 我们还要指出：进一步地，$a_n=0$ 对一切 $n\le0$（即包括 $n=0$）成立，当且仅当对一切 $N\in\R_{>0}$，
\[
|f(x+iy)|\le Cy^{-N}
\]
对某 $C\in\R_{>0}$ 与 $y\to\infty$ 成立。
\item 换言之，$f$ 是\textbf{速降}的\znote{「适度增长」与「速降」只差常数项 $a_0$ 是否为零：前者允许 $y$ 的多项式增长，后者要求对任意 $N$ 都有 $y^{-N}$ 级别的衰减。这正是第 3 讲中「尖点处全纯」与「尖点形式」的区别在分析层面的体现：去掉 $n\le0$ 的 Fourier 系数就得到速降。[来源：原文第 29 页；参见第 3 讲第 14、22 页]}。
\end{itemize}

\textbf{一般约化群的适度增长：}对 $\Q$ 上任意的约化群 $G$，适度增长条件可以这样表述。
\begin{itemize}
\item 对每个嵌入 $\rho:G\hookrightarrow\SL_{n,\Q}$（不是到 $\GL_{n,\Q}$），我们得到 $G(\A)$ 上的范数
\[
|g|:=\sup_v\max_{i,j=1,\dots,n}\{|\rho(g)_{i,j}|_v\},
\]
其中 $v$ 遍历一切素数与 $\infty$。
% ===== PDF p.030 =====
\item 函数 $\phi:G(\A)\to\C$ 称为\textbf{适度增长}（或缓慢增长）的，若
\[
|\phi(g)|\le C|g|^N
\]
对适当的 $C,N\in\R_{>0}$ 成立，参见 \texttt{[GH19, Definition 6.4.]}。
\item 不难看出，这个定义不依赖嵌入 $\rho$（因为换一个 $\rho$ 只会在相差某个多项式函数的意义下改变范数）。
\end{itemize}

我们把它留作练习（参见 \texttt{[GH19, Lemma 6.3.1.]} 以及上面关于 Fourier 系数的注记）：检验一个（全纯）截面
\[
f\in H^0(\GL_2(\Q)\backslash(\GL_2(\Af)\times\HH^\pm),\omega^{\otimes k})
\]
在尖点处全纯，当且仅当函数 $\phi_f\in C^\infty(\GL_2(\Q)\backslash\GL_2(\A))$ 具有适度增长。

我们简短总结一下已完成的工作。

\textbf{要点：}
\begin{itemize}
\item 我们可以描述
\[
M_k\to C^\infty(\GL_2(\Q)\backslash\GL_2(\A)),\qquad f\mapsto\phi_f
\]
的像。
\item 也就是说，$\phi\in C^\infty(\GL_2(\Q)\backslash\GL_2(\A))$ 落在像中，当且仅当
\begin{itemize}
\item $\phi(g,g_\infty z)=z^{-k}\phi(g,g_\infty)$ 对 $(g,g_\infty)\in\GL_2(\A)$、$z\in\C^\times$ 成立。
\item $\phi$ 具有适度增长。
\item 对每个 $g\in\GL_2(\Af)$，
\[
Y*\phi(g,-)=0,
\]
其中 $Y\in\g_\C=\mathfrak{gl}_2(\R)_\C$ 是前面构造的元素。
\end{itemize}
\item 注意，按我们的讨论，这些条件恰好对应模形式定义中的「模性、在尖点处全纯、全纯」。
\item 对 $k\ge1$，不难看出存在唯一的（在同构意义下）不可约 $(\g_\C,O_2(\R))$-模 $D'_{k-1}$，它含非零元素 $v$，使得 $Y*v=0$，且 $\C^\times$ 通过特征 $z\mapsto z^{-k}$ 作用在 $v$ 上，参见 \texttt{[Del73, Section 2.1]}（其中描述了 $D'_{k-1}$ 的对偶）\znote{$D'_{k-1}$ 是 $(\g_\C,O_2(\R))$-模：$O_2(\R)$ 给出 $K_\infty$-类型，$\g_\C$ 给出 Lie 代数作用；$Y*v=0$ 即最低权向量条件，$z\mapsto z^{-k}$ 则固定了权。于是上面的要点把「权 $k$ 的模形式」识别为从一个极小的表示 $D'_{k-1}$ 到适度增长函数空间的 $(\g_\C,O_2(\R))$-模同态，这正是把模形式纳入自守表示框架的关键一步；$D'_{k-1}$ 的构造在原文第 11 讲讨论。[来源：原文第 30 页]}。
\item 利用 $D'_{k-1}$，我们的结果可以改写为
\[
M_k\cong\Hom_{(\g_\C,O_2(\R))}\bigl(D'_{k-1},C^\infty_{\mathrm{mg}}(\GL_2(\Q)\backslash\GL_2(\A))\bigr),
\]
其中 $C^\infty_{\mathrm{mg}}(\GL_2(\Q)\backslash\GL_2(\A))\subseteq C^\infty(\GL_2(\Q)\backslash\GL_2(\A))$ 表示具有适度增长的函数构成的子集。
\item 更多内容将在我们讨论自守形式与 $(\g_\C,K_\infty)$-模时出现，参见第 11 讲。
\end{itemize}

我们以概述如何把尖点性纳入进来来结束本讲。

\textbf{尖点性的一个条件：}
\begin{itemize}
\item 设 $\phi=\phi_f\in C^\infty(\GL_2(\Q)\backslash\GL_2(\A))$，其中 $f\in M_k$，并把 $\phi$ 视为满足模性的函数 $\GL_2(\Af)\times\HH^\pm\to\C$。
\item 则 $f\in S_k$ 当且仅当积分
\[
\phi_B(g):=\int_{N(\Q)\backslash N(\A)}\phi(ng)\,dn
\]
% ===== PDF p.031 =====
对一切 $g\in\GL_2(\A)$ 为零；这里
\[
N=\begin{pmatrix}1&*\\0&1\end{pmatrix}
\]
是标准 Borel 子群
\[
B=\begin{pmatrix}*&*\\0&*\end{pmatrix}
\]
中的幺幂根基，$dn$ 是 $N(\Q)\backslash N(\A)$ 上的 $N(\A)$-不变测度\znote{把 $\phi$ 沿幺幂根基 $N$ 积分再令其为零，就是经典意义下「在所有尖点处常数 Fourier 系数为零」的表示论写法；下文的计算表明 $\phi_B(1)=\mu(m\hat\Z)a_0$，于是 $\phi_B\equiv0$ 恰好等价于各尖点处 $a_0=0$，即 $f$ 是尖点形式。[来源：原文第 30--31 页]}。
\item 也就是说：
\begin{itemize}
\item 可以假设 $f$（从而 $\phi=\phi_f$）在某个主同余子群 $K(m)\subseteq\GL_2(\Af)$ 下不变。
\item 取 $g=1$。利用 $N(\Q)\backslash N(\A)\cong(m\hat\Z\times\R)/m\Z$：
\begin{align*}
\int_{N(\Q)\backslash N(\A)}\phi_f(n)\,dn
&=\int_{N(\Q)\backslash N(\A)}j(n_\infty,i)^{-k}f(n,n_\infty i)\,d(n,n_\infty)\\
&=\int_{\R/m\Z}\int_{m\hat\Z}f(a,i+x)\,da\,dx\\
&=\mu(m\hat\Z)\int_{\R/m\Z}f(1,i+x)\,dx\\
&=\mu(m\hat\Z)a_0,
\end{align*}
其中 $\mu(m\hat\Z)$ 是 $m\hat\Z$ 的体积，而
\[
f(1,z)=\sum_{n=0}^\infty a_nq^{n/m}
\]
是模形式 $f(1,-)$ 关于 $\Gamma(m)$ 在尖点 $\infty$ 处的 Fourier 展开。
\item 确实：关键在于
\[
\int_{\R/m\Z}e^{2\pi i\frac nmx}\,dx=0
\]
当 $n>0$ 时成立（这是 Fourier 理论的一部分；更一般地，它可由「对任意紧 abelian 群 $A$ 与任意非平凡特征 $\chi:A\to\C^\times$，有 $\int_A\chi(g)\,dg=0$」推出）。
\item 于是条件「对一切 $g\in\GL_2(\A)$ 有 $\phi_B(g)=0$」等价于所有尖点处的常数 Fourier 系数消失。
\item 细节留作练习，参见 \texttt{[Del73, Rappel 1.3.2.]} 与 \texttt{[Gel75, Proposition 3.1.(vii)]}。
\end{itemize}
\item 注：因为 $\phi$ 对左平移是 $\GL_2(\Q)$-不变的，我们可以把 $B$ 换成任意共轭 $gBg^{-1}$，$g\in\GL_2(\Q)$。
\end{itemize}

我们要指出，我们把函数与 $\GL_2(\A)$ 相联系的方式在表述上不同于标准文献，例如 \texttt{[Gel75]}，后者通常限制到
\[
\Gamma_1(m):=\Bigl\{\begin{pmatrix}a&b\\c&d\end{pmatrix}\in\SL_2(\Z)\ \Bigm|\ c\equiv0,\ d\equiv1\ \mathrm{mod}\ m\Bigr\}
\]
% ===== PDF p.032 =====
上的模形式，或限制到关于
\[
\Gamma_0(m):=\Bigl\{\begin{pmatrix}a&b\\c&d\end{pmatrix}\in\SL_2(\Z)\ \Bigm|\ c\equiv0\ \mathrm{mod}\ m\Bigr\}
\]
带旁特征\idx{旁特征} $\chi:(\Z/m)^\times\to\C^\times$ 的模形式。尽管这些抓住了由模形式生成的一切自守表示，我们仍觉得在此刻把自己限制到这些同余子群（或它们在 $\GL_2(\Af)$ 中的对应物 $K_1(m),K_0(m)$）有点缺乏动机。下一讲我们将澄清我们的方法与经典方法之间的关系。
